\textbf{Phần 1. Trắc nghiệm lựa chọn}
\setcounterpageref{ex}{0}

\begin{ex} %[5D1V1]
Một số tự nhiên có $8$ chữ số. Chữ số hàng chục triệu là $6$, chữ số hàng triệu bé hơn chữ số hàng chục triệu $2$ đơn vị. Lớp nghìn gồm ba chữ số lẻ liên tiếp tăng dần từ hàng trăm nghìn đến hàng nghìn, biết chữ số hàng nghìn là $5$. Lớp đơn vị gồm toàn các chữ số $0$. Số đó là
\choice
{$64\,531\,000$}
{\True $64\,135\,000$}
{$62\,135\,000$}
{$64\,315\,000$}
\loigiai{
\begin{itemize}
    \item Lớp triệu: Chữ số hàng chục triệu là $6$. Chữ số hàng triệu là $6 - 2 = 4$.
    \item Lớp nghìn: Ba chữ số lẻ liên tiếp tăng dần kết thúc bằng chữ số hàng nghìn là $5$, do đó ba chữ số đó lần lượt là $1; 3; 5$ (hàng trăm nghìn là $1$, hàng chục nghìn là $3$, hàng nghìn là $5$).
    \item Lớp đơn vị: Gồm ba chữ số $0$.
\end{itemize}
Ghép lại ta được số: $64\,135\,000$.\\
\textbf{Đáp án của câu hỏi là B}
}
\end{ex}

\begin{ex} %[5D1V1]
Trong số tự nhiên $48\,725\,139$, giá trị của chữ số $8$ gấp giá trị của chữ số $2$ số lần là
\choice
{$40$ lần}
{$4\,000$ lần}
{\True $400$ lần}
{$40\,000$ lần}
\loigiai{
Trong số $48\,725\,139$:
\begin{itemize}
    \item Chữ số $8$ thuộc hàng triệu nên có giá trị là $8\,000\,000$.
    \item Chữ số $2$ thuộc hàng chục nghìn nên có giá trị là $20\,000$.
\end{itemize}
Giá trị của chữ số $8$ gấp giá trị của chữ số $2$ số lần là:
\[
8\,000\,000 : 20\,000 = 400\text{ (lần)}.
\]
\textbf{Đáp án của câu hỏi là C}
}
\end{ex}

\begin{ex} %[5D1V1]
Cho các thẻ số: $0$, $3$, $5$, $6$, $8$, $9$. Dùng tất cả $6$ thẻ số này để lập một số tự nhiên chẵn có $6$ chữ số khác nhau. Số lớn nhất có thể lập được là
\choice
{\True $986\,530$}
{$986\,503$}
{$985\,630$}
{$986\,350$}
\loigiai{
Để lập số có $6$ chữ số khác nhau lớn nhất:
\begin{itemize}
    \item Ta ưu tiên chọn các chữ số lớn nhất cho các hàng cao: chữ số hàng trăm nghìn là $9$, hàng chục nghìn là $8$, hàng nghìn là $6$, hàng trăm là $5$, hàng chục là $3$.
    \item Chữ số hàng đơn vị là $0$ (thỏa mãn số chẵn).
\end{itemize}
Vậy số chẵn lớn nhất có $6$ chữ số khác nhau lập được là $986\,530$.\\
\textbf{Đáp án của câu hỏi là A}
}
\end{ex}

\begin{ex} %[5D1V1]
Một số tự nhiên có bốn chữ số. Nếu viết thêm chữ số $7$ vào bên trái số đó thì ta được số mới lớn hơn số đã cho bao nhiêu đơn vị?
\choice
{$7\,000$}
{$700$}
{\True $70\,000$}
{$700\,000$}
\loigiai{
Gọi số tự nhiên có bốn chữ số ban đầu là $\overline{abcd}$.\\
Khi viết thêm chữ số $7$ vào bên trái số đó, ta được số mới có năm chữ số là $\overline{7abcd}$.\\
Phân tích cấu tạo số:
\[
\overline{7abcd} = 70\,000 + \overline{abcd}.
\]
Vậy số mới lớn hơn số ban đầu $70\,000$ đơn vị.\\
\textbf{Đáp án của câu hỏi là C}
}
\end{ex}

\begin{ex} %[5D1V1]
Một số tự nhiên có sáu chữ số có dạng $\overline{abc\,abc}$ (với $a 
e 0$). Khi phân tích theo cấu tạo số, số đó luôn bằng tích của số $\overline{abc}$ với số nào sau đây?
\choice
{$101$}
{$1\,010$}
{$10\,001$}
{\True $1\,001$}
\loigiai{
Phân tích cấu tạo số $\overline{abc\,abc}$:
\[
\overline{abc\,abc} = \overline{abc\,000} + \overline{abc} = \overline{abc} \times 1\,000 + \overline{abc} \times 1 = \overline{abc} \times (1\,000 + 1) = \overline{abc} \times 1\,001.
\]
Vậy số đó luôn bằng tích của $\overline{abc}$ nhân với $1\,001$.\\
\textbf{Đáp án của câu hỏi là D}
}
\end{ex}

\begin{ex} %[5D1V1]
Trên một đoạn phố, các ngôi nhà ở dãy bên phải được đánh số bằng các số chẵn liên tiếp, bắt đầu từ số nhà $12$ đến số nhà $108$. Dãy bên phải của đoạn phố đó có tất cả bao nhiêu ngôi nhà?
\choice
{$48$}
{\True $49$}
{$50$}
{$97$}
\loigiai{
Các số nhà ở dãy bên phải lập thành một dãy số chẵn cách đều $2$ đơn vị:
\[
12, 14, 16, \dots, 108.
\]
Số lượng ngôi nhà ở dãy phố đó là:
\[
(108 - 12) : 2 + 1 = 96 : 2 + 1 = 49\text{ (ngôi nhà)}.
\]
\textbf{Đáp án của câu hỏi là B}
}
\end{ex}

\begin{ex} %[5D1V1]
Một số tự nhiên khi làm tròn đến hàng chục nghìn thì được số $260\,000$. Số tự nhiên đó lớn nhất có thể là số nào sau đây?
\choice
{\True $264\,999$}
{$269\,999$}
{$264\,000$}
{$259\,999$}
\loigiai{
Khi làm tròn số đến hàng chục nghìn, ta xét chữ số hàng nghìn:
\begin{itemize}
    \item Nếu chữ số hàng nghìn bé hơn $5$ ($0, 1, 2, 3, 4$), ta làm tròn xuống bằng cách giữ nguyên chữ số hàng chục nghìn là $6$ và viết các chữ số sau thành $0$.
    \item Để số ban đầu có giá trị lớn nhất thì chữ số hàng chục nghìn là $6$, chữ số hàng nghìn lớn nhất có thể là $4$, và các chữ số ở các hàng còn lại (hàng trăm, chục, đơn vị) đều là $9$.
\end{itemize}
Vậy số tự nhiên lớn nhất thỏa mãn là $264\,999$.\\
\textbf{Đáp án của câu hỏi là A}
}
\end{ex}

\begin{ex} %[5D1V1]
Một cửa hàng điện máy bán một chiếc máy giặt giá $8\,650\,000$ đồng và một chiếc ti vi giá $14\,780\,000$ đồng. Bác Nam mua cả hai món đồ trên. Sau khi tính tổng số tiền, nhân viên thu ngân làm tròn số tiền đến hàng triệu. Số tiền sau khi làm tròn là
\choice
{$24\,000\,000$ đồng}
{\True $23\,000\,000$ đồng}
{$23\,500\,000$ đồng}
{$23\,400\,000$ đồng}
\loigiai{
Tổng số tiền thực tế mua máy giặt và ti vi là:
\[
8\,650\,000 + 14\,780\,000 = 23\,430\,000\text{ (đồng)}.
\]
Xét số $23\,430\,000$: chữ số hàng trăm nghìn là $4 < 5$, nên khi làm tròn đến hàng triệu ta giữ nguyên chữ số hàng triệu là $3$ và viết các chữ số phía sau thành $0$.\\
Vậy số tiền sau khi làm tròn là $23\,000\,000$ đồng.\\
\textbf{Đáp án của câu hỏi là B}
}
\end{ex}

\begin{ex} %[5D1V1]
Cho bốn số tự nhiên:
\begin{align*}
    M &= 40\,000\,000 + 500\,000 + 6\,000 + 20\\
    N &= 40\,560\,020\\
    P &= 40\,000\,000 + 50\,000 + 6\,000 + 2\\
    Q &= 40\,506\,200
\end{align*}
Số có giá trị lớn nhất trong bốn số trên là
\choice
{$M$}
{\True $N$}
{$P$}
{$Q$}
\loigiai{
Viết các số $M$ và $P$ dưới dạng số tự nhiên:
\begin{itemize}
    \item $M = 40\,506\,020$.
    \item $N = 40\,560\,020$.
    \item $P = 40\,056\,002$.
    \item $Q = 40\,506\,200$.
\end{itemize}
So sánh các số: $40\,056\,002 < 40\,506\,020 < 40\,506\,200 < 40\,560\,020$.\\
Vậy số lớn nhất là $N = 40\,560\,020$.\\
\textbf{Đáp án của câu hỏi là B}
}
\end{ex}

\begin{ex} %[5D1V1]
Để đánh số trang một cuốn truyện thiếu nhi dày $120$ trang (bắt đầu từ trang $1$ đến trang $120$), người ta phải dùng tất cả bao nhiêu lượt chữ số?
\choice
{$240$}
{$249$}
{\True $252$}
{$260$}
\loigiai{
Để đánh số trang từ trang $1$ đến trang $120$:
\begin{itemize}
    \item Từ trang $1$ đến trang $9$ có $9$ trang có $1$ chữ số, cần: $9 \times 1 = 9$ chữ số.
    \item Từ trang $10$ đến trang $99$ có $(99 - 10 + 1) = 90$ trang có $2$ chữ số, cần: $90 \times 2 = 180$ chữ số.
    \item Từ trang $100$ đến trang $120$ có $(120 - 100 + 1) = 21$ trang có $3$ chữ số, cần: $21 \times 3 = 63$ chữ số.
\end{itemize}
Tổng số lượt chữ số cần dùng là:
\[
9 + 180 + 63 = 252\text{ (chữ số)}.
\]
\textbf{Đáp án của câu hỏi là C}
}
\end{ex}

\begin{ex} %[5D1V1]
Cho số $749\,215$. Nếu chỉ được đổi chỗ đúng hai chữ số cho nhau thì số lớn nhất nhận được mà vẫn nhỏ hơn số ban đầu là số nào sau đây?
\choice
{\True $749\,125$}
{$729\,415$}
{$749\,152$}
{$719\,245$}
\loigiai{
Số ban đầu là $749\,215$. Để số mới nhận được nhỏ hơn số ban đầu nhưng lớn nhất trong các số nhỏ hơn:
\begin{itemize}
    \item Ta cần giữ nguyên các chữ số ở hàng cao nhất có thể (giữ hàng trăm nghìn là $7$, hàng chục nghìn là $4$, hàng nghìn là $9$).
    \item Xét ba chữ số còn lại ở lớp đơn vị: $2$, $1$, $5$. Đổi chỗ chữ số $2$ (hàng trăm) với chữ số $1$ (hàng chục) ta được số $749\,125$.
    \item Rõ ràng $749\,125 < 749\,215$ và lớn hơn các cách đổi chữ số ở hàng chục nghìn hoặc hàng trăm nghìn (ví dụ $729\,415$ hay $719\,245$).
\end{itemize}
Vậy số thỏa mãn là $749\,125$.\\
\textbf{Đáp án của câu hỏi là A}
}
\end{ex}

\begin{ex} %[5D1V1]
Bảng số liệu cho biết khoảng cách gần đúng từ Mặt Trời đến bốn hành tinh trong Hệ Mặt Trời:
\begin{center}
\begin{tabular}{|c|c|}
\hline
\textbf{Hành tinh} & \textbf{Khoảng cách đến Mặt Trời} \\
\hline
Sao Thủy & $57\,910\,000\text{ km}$ \\
\hline
Sao Kim & $108\,200\,000\text{ km}$ \\
\hline
Trái Đất & $149\,600\,000\text{ km}$ \\
\hline
Sao Hỏa & $227\,940\,000\text{ km}$ \\
\hline
\end{tabular}
\end{center}
Hành tinh nào có khoảng cách đến Mặt Trời gấp khoảng $2$ lần khoảng cách từ Sao Kim đến Mặt Trời?
\choice
{Trái Đất}
{\True Sao Hỏa}
{Sao Thủy}
{Sao Kim}
\loigiai{
Khoảng cách từ Sao Kim đến Mặt Trời là $108\,200\,000\text{ km}$.\\
Khoảng cách gấp $2$ lần khoảng cách đó là:
\[
108\,200\,000 \times 2 = 216\,400\,000\text{ (km)}.
\]
So sánh với các hành tinh còn lại: Khoảng cách từ Sao Hỏa đến Mặt Trời là $227\,940\,000\text{ km}$, xấp xỉ gần nhất với $216\,400\,000\text{ km}$ (gấp khoảng $2$ lần khoảng cách từ Sao Kim).\\
\textbf{Đáp án của câu hỏi là B}
}
\end{ex}

\textbf{Phần 2. Trắc nghiệm đúng sai}
\setcounterpageref{ex}{0}

\begin{ex} %[5D1V1]
Cho năm thẻ số: $0$, $2$, $5$, $7$, $8$. Người ta dùng cả năm thẻ số này để lập các số tự nhiên có năm chữ số khác nhau. Xét tính đúng sai của các khẳng định sau:
\choiceTF
{\True Số lẻ lớn nhất có thể lập được là $87\,205$.}
{\True Số bé nhất chia hết cho $2$ có thể lập được là $20\,578$.}
{\True Có tất cả $24$ số có $5$ chữ số khác nhau có chữ số tận cùng bằng $0$.}
{Hiệu giữa số chẵn lớn nhất và số chẵn bé nhất lập được là $67\,230$.}
\loigiai{
\begin{itemize}
    \item a) Để lập số lẻ lớn nhất có $5$ chữ số khác nhau, chữ số tận cùng phải là chữ số lẻ ($5$ hoặc $7$). Để số là lớn nhất, ta chọn chữ số hàng chục nghìn là $8$, hàng nghìn là $7$, hàng trăm là $2$, hàng chục là $0$ và hàng đơn vị là $5$. Ta được số $87\,205$. Khẳng định này \textbf{Đúng}.
    \item b) Số chia hết cho $2$ phải có chữ số tận cùng là $0, 2$ hoặc $8$. Để số là bé nhất, chữ số hàng chục nghìn phải bé nhất khác $0$, tức là chọn chữ số $2$. Khi đó chữ số tận cùng có thể là $0$ hoặc $8$. Nếu tận cùng là $8$, số bé nhất lập được là $20\,578$. Nếu tận cùng là $0$, số bé nhất là $25\,780 > 20\,578$. Vậy số bé nhất chia hết cho $2$ là $20\,578$. Khẳng định này \textbf{Đúng}.
    \item c) Các số có tận cùng bằng $0$ có dạng $\overline{abcd0}$. Bốn chữ số $a, b, c, d$ được chọn từ bốn chữ số $\{2, 5, 7, 8\}$: hàng chục nghìn có $4$ cách chọn, hàng nghìn có $3$ cách chọn, hàng trăm có $2$ cách chọn, hàng chục có $1$ cách chọn. Số các số lập được là $4 \times 3 \times 2 \times 1 = 24$ (số). Khẳng định này \textbf{Đúng}.
    \item d) Số chẵn lớn nhất lập được là $87\,520$. Số chẵn bé nhất lập được là $20\,578$. Hiệu giữa hai số là:
    \[
    87\,520 - 20\,578 = 66\,942 
e 67\,230.
    \]
    Khẳng định này \textbf{Sai}.
\end{itemize}
\begin{itemize}
    \item a) \textbf{Đúng}.
    \item b) \textbf{Đúng}.
    \item c) \textbf{Đúng}.
    \item d) \textbf{Sai}: Hiệu đúng là $66\,942$.
\end{itemize}
}
\end{ex}

\begin{ex} %[5D1V1]
Sản lượng cà phê xuất khẩu của một doanh nghiệp trong ba quý đầu năm được ghi nhận như sau: Quý 1 xuất khẩu $125\,460$ tấn; Quý 2 xuất khẩu $148\,720$ tấn; Quý 3 xuất khẩu $132\,890$ tấn. Xét tính đúng sai của các khẳng định sau:
\choiceTF
{\True Sản lượng Quý 2 nhiều hơn sản lượng Quý 1 là $23\,260$ tấn.}
{Làm tròn sản lượng Quý 1 đến hàng chục nghìn ta được $120\,000$ tấn.}
{\True Làm tròn sản lượng Quý 2 đến hàng trăm nghìn ta được $100\,000$ tấn.}
{\True Nếu làm tròn sản lượng mỗi quý đến hàng nghìn rồi tính tổng thì ước tính cả ba quý doanh nghiệp xuất khẩu được $407\,000$ tấn cà phê.}
\loigiai{
\begin{itemize}
    \item Hiệu sản lượng Quý 2 và Quý 1: $148\,720 - 125\,460 = 23\,260$ (tấn). Khẳng định a \textbf{Đúng}.
    \item Số $125\,460$ có chữ số hàng nghìn là $5 \ge 5$, nên khi làm tròn đến hàng chục nghìn ta được số $130\,000$ tấn. Khẳng định b \textbf{Sai}.
    \item Số $148\,720$ có chữ số hàng chục nghìn là $4 < 5$, nên khi làm tròn đến hàng trăm nghìn ta được số $100\,000$ tấn. Khẳng định c \textbf{Đúng}.
    \item Làm tròn sản lượng từng quý đến hàng nghìn:
    \begin{itemize}
        \item Quý 1: $125\,460$ làm tròn thành $125\,000$ tấn (chữ số hàng trăm là $4 < 5$).
        \item Quý 2: $148\,720$ làm tròn thành $149\,000$ tấn (chữ số hàng trăm là $7 \ge 5$).
        \item Quý 3: $132\,890$ làm tròn thành $133\,000$ tấn (chữ số hàng trăm là $8 \ge 5$).
    \end{itemize}
    Tổng sản lượng ước tính là:
    \[
    125\,000 + 149\,000 + 133\,000 = 407\,000\text{ (tấn)}.
    \]
    Khẳng định d \textbf{Đúng}.
\end{itemize}
\begin{itemize}
    \item a) \textbf{Đúng}.
    \item b) \textbf{Sai}: Làm tròn đến hàng chục nghìn được $130\,000$ tấn.
    \item c) \textbf{Đúng}.
    \item d) \textbf{Đúng}.
\end{itemize}
}
\end{ex}

\textbf{Phần 3. Trả lời ngắn}
\setcounterpageref{ex}{0}

\begin{ex} %[5D1V1]
Tìm số tự nhiên bé nhất có các chữ số khác nhau mà tổng các chữ số của số đó bằng $25$.
\shortans{1789}
\loigiai{
\textbf{Đáp số:} $1789$\\
Để tìm số tự nhiên bé nhất có tổng các chữ số bằng $25$:
\begin{itemize}
    \item Số đó phải có ít chữ số nhất có thể.
    \item Để có ít chữ số nhất thì các chữ số ở hàng thấp phải lớn nhất có thể: $9 + 8 + 7 = 24$.
    \item Chữ số hàng cao nhất còn lại là: $25 - 24 = 1$.
\end{itemize}
Sắp xếp các chữ số theo thứ tự từ bé đến lớn để được số bé nhất: $1\,789$.
}
\end{ex}

\begin{ex} %[5D1V1]
Một số tự nhiên khi làm tròn đến hàng trăm nghìn thì được $700\,000$. Biết số đó là số lẻ bé nhất có thể. Hỏi số tự nhiên ban đầu là số nào?
\shortans{650001}
\loigiai{
\textbf{Đáp số:} $650001$\\
Để một số tự nhiên khi làm tròn đến hàng trăm nghìn được $700\,000$ thì:
\begin{itemize}
    \item Chữ số hàng trăm nghìn có thể là $6$ (làm tròn lên) hoặc $7$ (làm tròn xuống).
    \item Để số là bé nhất có thể, số đó phải có chữ số hàng trăm nghìn là $6$ và chữ số hàng chục nghìn nhỏ nhất phải là $5$ (để được làm tròn lên thành $7$).
    \item Các chữ số tiếp theo ở hàng nghìn, hàng trăm, hàng chục là $0$. Để số là số lẻ bé nhất thì chữ số hàng đơn vị phải là $1$.
\end{itemize}
Vậy số cần tìm là $650\,001$.
}
\end{ex}

\begin{ex} %[5D1V1]
Trong các số tự nhiên có ba chữ số từ $100$ đến $999$, có tất cả bao nhiêu số có chữ số hàng trăm là chữ số chẵn?
\shortans{400}
\loigiai{
\textbf{Đáp số:} $400$\\
Gọi số có ba chữ số là $\overline{abc}$ ($100 \le \overline{abc} \le 999$):
\begin{itemize}
    \item Chữ số hàng trăm $a$ là chữ số chẵn khác $0$, nên $a \in \{2, 4, 6, 8\}$ ($4$ cách chọn).
    \item Chữ số hàng chục $b \in \{0, 1, 2, \dots, 9\}$ ($10$ cách chọn).
    \item Chữ số hàng đơn vị $c \in \{0, 1, 2, \dots, 9\}$ ($10$ cách chọn).
\end{itemize}
Số các số tự nhiên thỏa mãn yêu cầu là:
\[
4 \times 10 \times 10 = 400\text{ (số)}.
\]
}
\end{ex}

\begin{ex} %[5D1V1]
Người ta dùng tất cả $282$ lượt chữ số để đánh số trang của một cuốn sách (bắt đầu từ trang $1$). Hỏi cuốn sách đó dày bao nhiêu trang?
\shortans{130}
\loigiai{
\textbf{Đáp số:} $130$\\
\begin{itemize}
    \item Từ trang $1$ đến trang $9$ có $9$ trang có $1$ chữ số, cần: $9 \times 1 = 9$ chữ số.
    \item Từ trang $10$ đến trang $99$ có $(99 - 10 + 1) = 90$ trang có $2$ chữ số, cần: $90 \times 2 = 180$ chữ số.
    \item Số chữ số còn lại để đánh các trang có $3$ chữ số là:
    \[
    282 - (9 + 180) = 93\text{ (chữ số)}.
    \]
    \item Số trang sách có $3$ chữ số là:
    \[
    93 : 3 = 31\text{ (trang)}.
    \]
\end{itemize}
Vậy cuốn sách đó dày số trang là: $99 + 31 = 130$ trang.
}
\end{ex}

\begin{ex} %[5D1V1]
Cho bốn thẻ số: $0$, $2$, $5$, $8$. Hỏi có thể lập được tất cả bao nhiêu số tự nhiên có ba chữ số khác nhau và là số chẵn?
\shortans{14}
\loigiai{
\textbf{Đáp số:} $14$\\
Gọi số chẵn có ba chữ số khác nhau là $\overline{abc}$ ($a 
e 0$, $c \in \{0, 2, 8\}$):
\begin{itemize}
    \item \textbf{Trường hợp 1:} $c = 0$.
    Chữ số $a$ có $3$ cách chọn (từ $\{2, 5, 8\}$). Chữ số $b$ có $2$ cách chọn từ hai chữ số còn lại. Số các số lập được là: $3 \times 2 = 6$ số.
    \item \textbf{Trường hợp 2:} $c = 2$.
    Chữ số $a 
e 0$ và $a 
e 2$ nên $a \in \{5, 8\}$ ($2$ cách chọn). Chữ số $b$ có $2$ cách chọn từ hai chữ số còn lại (kể cả số $0$). Số các số lập được là: $2 \times 2 = 4$ số.
    \item \textbf{Trường hợp 3:} $c = 8$.
    Chữ số $a 
e 0$ và $a 
e 8$ nên $a \in \{2, 5\}$ ($2$ cách chọn). Chữ số $b$ có $2$ cách chọn từ hai chữ số còn lại. Số các số lập được là: $2 \times 2 = 4$ số.
\end{itemize}
Tổng số các số chẵn lập được là:
\[
6 + 4 + 4 = 14\text{ (số)}.
\]
}
\end{ex}

\begin{ex} %[5D1V1]
Một trường tiểu học dự kiến mua $45$ bộ bàn ghế học sinh mới với đơn giá $380\,000$ đồng một bộ. Kế toán của trường làm tròn tổng số tiền cần thanh toán đến hàng triệu để lập dự trù kinh phí chuyển khoản. Hỏi số tiền sau khi làm tròn là bao nhiêu đồng?
\shortans{17000000}
\loigiai{
\textbf{Đáp số:} $17000000$\\
Tổng số tiền thực tế để mua $45$ bộ bàn ghế là:
\[
45 \times 380\,000 = 17\,100\,000\text{ (đồng)}.
\]
Xét số $17\,100\,000$: chữ số hàng trăm nghìn là $1 < 5$, do đó khi làm tròn đến hàng triệu ta giữ nguyên chữ số hàng triệu là $7$ và viết các chữ số phía sau thành $0$.\\
Vậy số tiền sau khi làm tròn là $17\,000\,000$ đồng.
}
\end{ex}
